\subsubsection{证书优化器与共终细化}\label{subsubsec:typed-address-biaxial-completion-certificate-optimizer-cofinal-refinement}
局部小值弧现在拥有多种可验证上界：单中心上界、packet 的 Hölder 上界、split 解析上界以及有限弧求积上界。优化层的任务不是提供额外数学不等式，而是在这些已验证局部证书之间选择 winner，并以严格区间算术闭合 Jensen 排斥阈值。

固定 \(0<R<r<1\)。令
\[
\mathcal T_{r,R}:=\log\frac rR,
\]
并把所有可调选择记为
\[
\Xi=
\left(
\mathcal P_\theta,\,
\mathcal C_{\mathrm{packet}},\,
\mathcal S_{\mathrm{cell}},\,
\mathcal M_{\mathrm{method}},\,
\mathcal T_{\mathrm{quad}},\,
p_{\mathrm{Taylor}},\,
b_{\mathrm{round}},\,
\mathcal Z_{\mathrm{box}}
\right).
\]
这些分量分别表示圆周弧划分、packet 聚类、cell 切分、局部上界方法、求积子区间、Taylor 阶数、区间算术位数以及零点参数盒。对应的全局上界写为
\[
\mathcal U(\Xi)
:=
J_r^+
+
\log\frac{A_0^+}{A_0^-}
+
B_{\mathrm{ZA}}^+(\Xi).
\]
证书裕度必须以区间端点核验：
\[
\mathfrak m(\Xi)
:=
\left[\mathcal T_{r,R}\right]^-
-
\left[\mathcal U(\Xi)\right]^+.
\]
只有
\[
\mathfrak m(\Xi)>0
\]
才允许输出通过。

\begin{definition}[优化证书]\label{def:typed-address-biaxial-completion-optimization-certificate}
一份优化证书记为
\[
\mathsf{OptCert}_{r,R}
:=
\left(
r,R,\delta,
\mathcal P_\theta,
\mathcal L_{\mathrm{local}},
B_{\mathrm{ZA,opt}}^+,
\mathcal U^+,
\mathfrak m,
\mathsf{TranscriptHash}
\right),
\]
其中
\[
\mathcal L_{\mathrm{local}}
:=
\{(I,\mathsf{LocalCert}_I,B_I^+):I\in\mathcal I_{\mathrm{bad}}\}.
\]
每个局部证书属于
\[
\mathsf{LocalCert}_I
\in
\{
\mathsf{SingleZeroCert},
\mathsf{PacketCert},
\mathsf{SplitCert},
\mathsf{IQCert}
\}.
\]
若同一弧 \(I\) 存在多个已验证候选，则最终 ledger 只使用
\[
B_I^+
=
\min
\left\{
B_{I,\alpha}^+:
\mathsf{LocalCert}_{I,\alpha}\text{ 已通过验证}
\right\}.
\]
\end{definition}

优化器可以使用启发式评分、随机搜索或非严格预测模型；这些选择不进入数学正确性。确定性验证器
\[
\mathsf{VerifyOpt}(\mathsf{OptCert}_{r,R})
\in
\{\mathrm{PASS},\mathrm{FAIL}\}
\]
只检查有限数据：局部证书逐个通过、弧覆盖正确、每个 \(B_I^+\) 是对应局部证书的严格上界、总和的区间算术正确，以及
\[
\mathcal U^+<\left[\log\frac rR\right]^-.
\]

\begin{theorem}[优化证书的零点排斥语义]\label{thm:typed-address-biaxial-completion-optimized-jensen-rejection-certificate}
若
\[
\mathsf{VerifyOpt}(\mathsf{OptCert}_{r,R})=\mathrm{PASS},
\]
则
\[
\Phi_a(w)\ne0
\qquad(|w|\le R).
\]
\end{theorem}
\begin{proof}
每个局部 winner 给出
\[
B_I\le B_I^+.
\]
求和得到
\[
B_{\mathrm{ZA}}\le B_{\mathrm{ZA,opt}}^+.
\]
验证器又确认
\[
J_r^+
+
\log\frac{A_0^+}{A_0^-}
+
B_{\mathrm{ZA,opt}}^+
<
\log\frac rR
\]
以严格区间算术成立。于是前述 zero-aware Jensen 排斥链给出 \(|w|\le R\) 内无零点。证毕。
\end{proof}

\begin{definition}[细化偏序]\label{def:typed-address-biaxial-completion-refinement-preorder}
若 \(\Xi'\) 相比 \(\Xi\) 只进行以下操作，则记为 \(\Xi'\preceq\Xi\)：弧划分细化；零点盒收缩；packet 被拆分或保留；cell 被二分或保留；求积子区间被二分或 Taylor 阶数增加；区间算术精度增加；局部模型加强，即
\[
c_K'\ge c_K,
\qquad
d'\le d.
\]
\end{definition}

单次细化后的数值上界未必单调下降。区间自动微分可能在某些小 cell 上更宽，高阶 Taylor 项也可能携带更大的舍入包络。因此优化层必须保留 incumbent ledger：
\[
B_I^{\mathrm{best}}(t+1)
:=
\min\left\{
B_I^{\mathrm{best}}(t),
B_{I}^{\mathrm{cand}}(t+1)
\right\}.
\]
于是
\[
B_{\mathrm{ZA}}^{\mathrm{best}}(t+1)
\le
B_{\mathrm{ZA}}^{\mathrm{best}}(t).
\]
该单调性来自历史最优保留，而不是来自任一细化动作本身。

同一弧上的候选还可按支配关系压缩。若两个已验证局部证书 \(C_1,C_2\) 满足
\[
B^+(C_1)\le B^+(C_2),
\qquad
\mathrm{size}(C_1)\le\mathrm{size}(C_2),
\]
且 \(C_1\) 的依赖数据不强于 \(C_2\)，则 \(C_2\) 可从候选池删除。未被支配的局部候选构成 Pareto frontier：
\[
\mathcal F_I
:=
\{C:C\text{ 未被其他已验证局部证书支配}\}.
\]
最终只需保留
\[
\min_{C\in\mathcal F_I}B^+(C).
\]

优化动作分为四类：
\[
\mathsf{Act}
=
\mathsf{Act}_{\mathrm{geo}}
\cup
\mathsf{Act}_{\mathrm{zero}}
\cup
\mathsf{Act}_{\mathrm{quad}}
\cup
\mathsf{Act}_{\mathrm{arith}}.
\]
几何动作包括弧二分、packet 拆分和 cell 二分，目标是提高 inactive shield 下界；零点盒动作收缩 \(\Theta_\ell,\Gamma_\ell\)，目标是降低 log budget；求积动作二分 \(J_j\) 或提高 Taylor 阶数，目标是降低余项
\[
\frac{2M_{j,p+1}h_j^{p+2}}{(p+1)!(p+2)};
\]
算术动作提高 \(b_{\mathrm{round}}\)，用于处理舍入导致的 strict margin 损失。

预算分配给出调度压力。定义可用于小值弧的局部预算
\[
T_{\mathrm{loc}}
:=
\log\frac rR
-
J_r^+
-
\log\frac{A_0^+}{A_0^-}
-
\frac12m(E_0)\log\left(1+\frac{\delta}{m_0}\right).
\]
取配额 \(\tau_I>0\)，满足
\[
\sum_I\tau_I=T_{\mathrm{loc}}.
\]
可按弧长分配，也可按预测真实贡献 \(\widehat B_I\) 分配。局部超额与压力为
\[
s_I:=B_I^{\mathrm{best}}-\tau_I,
\]
\[
\mathsf{Press}(I)
:=
s_I^+
+
\alpha_{\mathrm{err}}E_I
+
\alpha_{\mathrm{box}}W_I
+
\alpha_{\mathrm{round}}R_I.
\]
这里 \(E_I\) 表示求积余项或区间宽度，\(W_I\) 表示零点盒宽度导致的对数预算过剩，\(R_I\) 表示舍入宽度。调度器优先处理压力最大的弧；该排序只影响搜索效率，不影响证明。

对动作 \(a\) 可定义预测单位成本收益
\[
\rho(a)
:=
\frac{\widehat\Delta(a)}
{\mathrm{cost}(a)+\varepsilon},
\]
其中
\[
\mathrm{cost}(a)
=
c_1\Delta N_{\mathrm{cell}}
+
c_2\Delta N_{\mathrm{AD}}
+
c_3\Delta p
+
c_4\Delta b_{\mathrm{round}}.
\]
\(\rho(a)\) 只用于调度；最终证书完全由已验证数据承担。

\begin{definition}[共终搜索族]\label{def:typed-address-biaxial-completion-cofinal-search-family}
令
\[
\mathfrak C
:=
\bigcup_{N,p,b}\mathfrak C(N,p,b),
\]
其中 \(N\) 控制最大 cell 数，\(p\) 控制最大 Taylor 阶数，\(b\) 控制区间算术精度。若调度满足
\[
\forall C\in\mathfrak C,\qquad C\text{ eventually generated},
\]
则称该调度对 \(\mathfrak C\) 共终。
\end{definition}

一个可实现的调度是按总等级
\[
N+p+\log_2 b
\]
做 dovetailing，同时在每个等级内按压力和预测收益排序。实际搜索可采用 best-first heuristic 加 cofinal fairness guard：大多数步骤追逐当前收益，周期性强制推进低等级尚未访问的 refinement。

\begin{theorem}[共终优化器的有限完备性]\label{thm:typed-address-biaxial-completion-cofinal-optimizer-completeness}
假设：
\begin{enumerate}
  \item 局部证书生成族对有限弧划分、packet 划分、cell splitting、Taylor degree 与 precision 共终；
  \item 所有局部 verifier 均 sound；
  \item 存在某个有限证书 \(\Xi_\star\)，其严格裕度满足 \(\mathfrak m_\star>0\)；
  \item 提高区间算术精度可使初等函数 enclosure 宽度趋于 \(0\)；
  \item 零点盒细化可达到 \(\Xi_\star\) 所需精度；
  \item 对 peeled smooth remainder 的区间 Taylor 求积上界收敛到真实积分。
\end{enumerate}
则任一 fair cofinal optimizer 在有限步内输出 \(\mathrm{PASS}\)。
\end{theorem}
\begin{proof}
由 \(\mathfrak m_\star>0\)，取 \(0<\eta<\mathfrak m_\star\)。只要所有局部上界与基点、阈值等区间量的总误差小于 \(\eta\)，严格不等式仍保持。共终性保证某个时刻会生成不粗于 \(\Xi_\star\) 的组合结构。零点盒细化、区间算术收敛与 Taylor 求积收敛保证对应数值包络进入 \(\eta\)-邻域。于是某一步有 \(\mathfrak m>0\)，验证器输出 \(\mathrm{PASS}\)。证毕。
\end{proof}

资源受限时，优化器输出 failure frontier，而非数学否证。定义
\[
\mathsf{OptFail}_{r,R}
:=
\left(
\mathrm{Caps},
\mathcal F_{\mathrm{frontier}},
\Omega,
\mathcal B_{\mathrm{bottleneck}},
\mathcal A_{\mathrm{blocked}}
\right),
\]
其中
\[
\Omega:=\mathcal U^+-\left[\mathcal T_{r,R}\right]^-.
\]
\(\mathcal F_{\mathrm{frontier}}\) 记录每个弧当前最优局部证书与 failure mode；
\[
\mathcal B_{\mathrm{bottleneck}}
:=
\{I:s_I>0\}
\]
记录瓶颈弧；\(\mathcal A_{\mathrm{blocked}}\) 记录被上限阻止的动作，如 \(p=p_{\max}\)、\(N_{\mathrm{cell}}=N_{\max}\) 或 \(b=b_{\max}\)。若尚未穷举 caps 内全部候选，则语义为 \(\mathrm{FrontierFail}\)；若指定 caps 内所有候选均已枚举且未通过，则语义为 \(\mathrm{CappedExhaustiveFail}\)，只能说明 caps 内不存在通过证书。

优化层失败类型可细分为
\[
\mathsf{OptFailType}
\in
\{
\mathrm{GlobalMarginFail},
\mathrm{LocalBudgetFail},
\mathrm{ZeroBoxFail},
\mathrm{PacketGeometryFail},
\mathrm{QuadratureFail},
\mathrm{RoundoffFail},
\mathrm{CertificateSizeFail},
\mathrm{ExhaustedCapsFail}
\}.
\]
\(\mathrm{GlobalMarginFail}\) 表示
\[
\log\frac rR
-
J_r^+
-
\log\frac{A_0^+}{A_0^-}
\le0.
\]
\(\mathrm{LocalBudgetFail}\) 记录 \(B_I^{\mathrm{best}}>\tau_I\) 的瓶颈弧。\(\mathrm{ZeroBoxFail}\) 指向 \(\Theta_\ell,\Gamma_\ell\) 宽度造成的 log budget 过剩。\(\mathrm{PacketGeometryFail}\) 指向过弱的 inactive shield 或过小的 \(C_K\)。\(\mathrm{QuadratureFail}\) 记录最大 \(E_j\) 的求积子区间。\(\mathrm{RoundoffFail}\) 记录 \(b_{\mathrm{round}}\)、\(\operatorname{wid}(\mathcal U)\) 与 \(\mathfrak m\)。\(\mathrm{CertificateSizeFail}\) 记录 size ledger。

Branch-and-bound 只作为效率工具。若对某搜索分支可证明乐观下界
\[
\underline B_{\mathrm{branch}}
\le
\inf_{\Xi\in\mathrm{branch}}B_{\mathrm{ZA}}^+(\Xi),
\]
并且
\[
J_r^+
+
\log\frac{A_0^+}{A_0^-}
+
\underline B_{\mathrm{branch}}
\ge
\log\frac rR,
\]
则该分支不可能通过，可以剪枝。这里必须使用下界而非上界；剪枝不是正确性必需项。

Packet clustering 也应共终枚举。若一个弧内零点按角度排序为 \(z_1,\ldots,z_N\)，任一 packet clustering 是区间划分
\[
[1,a_1]\sqcup[a_1+1,a_2]\sqcup\cdots\sqcup[a_s+1,N].
\]
共有 \(2^{N-1}\) 种。启发式可按 gap score
\[
g_j
:=
\frac{\operatorname{dist}(\Theta_j,\Theta_{j+1})}
{\operatorname{diam}\Theta_j+\operatorname{diam}\Theta_{j+1}
+\Gamma_j^++\Gamma_{j+1}^+}
\]
优先切开大 gap，但 cofinal guard 必须保证所有有序 packet 划分最终被尝试。

方法切换的安全规则为
\[
B_K^{\mathrm{best}}
:=
\min
\left\{
B_{K,\mathrm{Holder}}^+,\,
B_{K,\mathrm{split}}^+,\,
B_{K,\mathrm{IQ}}^+
\right\}_{\mathrm{verified}}.
\]
未通过的方法只留下调度信息，不进入证明 ledger。最终证书只需输出每个弧的 winner
\[
C_I^\star:=\arg\min_jB^+(C_{I,j})
\]
以及 transcript hash；大量搜索候选不必进入数学证明核心。

所有标量必须以区间形式输出。全局通过条件采用
\[
\left[
J_r^+
+
\log\frac{A_0^+}{A_0^-}
+
B_{\mathrm{ZA,opt}}^+
\right]^+
<
\left[\log\frac rR\right]^-.
\]
定义 roundoff slack
\[
\mathfrak r
:=
\left[\log\frac rR\right]^-
-
\left[
J_r^+
+
\log\frac{A_0^+}{A_0^-}
+
B_{\mathrm{ZA,opt}}^+
\right]^+.
\]
必须有 \(\mathfrak r>0\)。若中点裕度为正但 \(\mathfrak r\le0\)，输出 \(\mathrm{RoundoffFail}\)，而非通过。

优化层的最终接入为
\[
B_{\mathrm{ZA,opt}}^+
:=
\frac12m(E_0)\log\left(1+\frac{\delta}{m_0}\right)
+
\sum_I B_I^{\mathrm{best}}.
\]
于是整体链条为
\[
\mathsf{RTCert}
\Longrightarrow
\mathsf{PACert}
\Longrightarrow
\mathsf{LocalCert}^{\star}
\Longrightarrow
\mathsf{OptCert}
\Longrightarrow
\mathsf{ZACert}_{\mathrm{opt}}
\Longrightarrow
\text{Jensen 排斥证书}.
\]
